% ===== 原文 pp.21–23 =====
% tr03b：承 tr03a。p.21 末三行（定理 9 的证明开头，仍在 p.21 上）置于本片开头，
% 使 amsthm 的 proof 环境（p.21→p.22）不跨文件；随后为 pp.22–23 正文。
% [原文 p.21]

\begin{proof}
对应于剪开曲面 $R(\gamma)$ 的各分支 $R'$，考虑映射类群的短正合列
\[
1\longrightarrow\prod_j\operatorname{Dehn}(\gamma_j)\longrightarrow\bigcap_j\Stab_0(\gamma_j)
\longrightarrow\prod_{R(\gamma)\ \text{components}}\MCG(R')\longrightarrow1,
\]
% [原文 p.22]
以及由 Fenchel--Nielsen 坐标给出的 Teichmüller 空间的相应纤维化
\[
\begin{array}{ccc}
\prod_{R(\gamma)\ \text{components}}\cT(R') & \hookrightarrow & \cT(R)\\[-2pt]
& & \big\downarrow\\[-8pt]
& & \prod_{\gamma_j}\R_{>0}\times\R.
\end{array}
\]
（该短正合列把半 Dehn 扭转放进单边界环面的映射类群中。）$\dd\ell\wedge\dd\tau$ 公式保证了该纤维化是辛流形之间的纤维化。

为建立该公式，先对陪集和写出
\[
\sum_{\MCG/\cap_j\Stab_0(\gamma_j)}f=\sum_{\MCG/\Stab(\gamma)}\ \sum_{\Stab(\gamma)/\cap_j\Stab_0(\gamma_j)}f
=|\operatorname{Sym}(\gamma)|\,f_\gamma,
\]
第二个等号用到 $f_\gamma$ 关于 $\operatorname{Sym}(\gamma)$ 不变。把所得的 $f_\gamma$ 表达式代入积分，并展开求和（即把 $\MCG/\cap_j\Stab_0(\gamma_j)$ 的平移和写成某个 MCG 基本域的各平移之和），得到等式
\[
\int_{\cT(R)/\MCG}f_\gamma\,\dd V
=(|\operatorname{Sym}(\gamma)|)^{-1}\int_{\cT(R)/\cap_j\Stab_0(\gamma_j)}f\,\dd V.
\]
代入纤维化
\[
\begin{array}{ccc}
\prod_{R(\gamma)\ \text{components}}\cT(R')/\MCG(R') & \hookrightarrow & \cT(R)/\bigcap_j\Stab_0(\gamma_j)\\[-2pt]
& & \big\downarrow\\[-8pt]
& & \prod_{\gamma_j}(\R_{>0}\times\R)/\operatorname{Dehn}_*(\gamma_j),
\end{array}
\]
其中若该曲线围出一个单边界环面，则 $\operatorname{Dehn}_*(\gamma_j)$ 由半扭转生成，否则由简单扭转生成。再代入体积元的分解
\[
\dd V=\prod_{R(\gamma)\ \text{components}}\dd V(R')\times\prod_{\gamma_j}\dd\ell_j\wedge\dd\tau_j.
\]
函数 $f$ 只依赖于各值 $\mathbf{x}$。固定这些值 $\mathbf{x}$，作 $\prod\cT(R')/\MCG(R')$ 上的积分，即得到乘积体积 $V(R(\gamma);\mathbf{x})$。最后，$\operatorname{Dehn}_*(\gamma_j)$ 只作用在变量 $\tau_j$ 上：若 $\gamma_j$ 围出单边界环面，则基本域为 $0<\tau_j<\ell_j/2$，否则基本域为 $0<\tau_j<\ell_j$。对单边界环面，这一作用通过改用"原 $V(L)$ 的一半"作为体积值来计入。公式的右端至此建立。
\end{proof}

% [Tutor Note: 这段证明有两个"不变量"技巧值得留意：(1) 先用陪集和把
%  $\MCG/\bigcap_j\Stab_0(\gamma_j)$ 的平移和拆成 $\MCG/\Stab(\gamma)$ 与
%  $\Stab(\gamma)/\bigcap_j\Stab_0(\gamma_j)$ 两层，第二层恰好有
%  $|\operatorname{Sym}(\gamma)|$ 个元素，于是前者的和等于 $|\operatorname{Sym}(\gamma)|f_\gamma$；
%  (2) 展开（unfold）求和之后，被积函数变成 $\cT(R)/\bigcap_j\Stab_0(\gamma_j)$ 上的单个积分，
%  再沿纤维化先对纤维积分。纤维基变量是 $(\ell_j,\tau_j)$：普通曲线的 $\operatorname{Dehn}_*$
%  由 $(\ell,\tau)\mapsto(\ell,\tau+\ell)$ 生成，基本域 $0<\tau<\ell$，$\tau$ 方向贡献因子 $\ell$；
%  围出单边界环面的 $\gamma_j$ 由半扭转 $(\ell,\tau)\mapsto(\ell,\tau+\ell/2)$ 生成，基本域宽
%  $\ell/2$，贡献因子 $\ell/2$——这就是 p.22 末"体积取 $1/2$"的几何来源，也是体积元分解
%  $\dd V=\prod\dd V(R')\times\prod\dd\ell_j\wedge\dd\tau_j$ 能逐项相乘的原因（由定理 2 的
%  $\dd\ell\wedge\dd\tau$ 公式，WP 形式在剪开后就是各部分形式之和）。]

% [原文 p.23]
现在可以把余体积公式应用到长度恒等式上。这一应用与亏格一的情形类似。同样地——关键在于使用 MCG 对测地长度函数的作用 $\ell_\alpha\circ h^{-1}=\ell_{h(\alpha)}$，从而把关于拓扑构型的求和看成函数的 MCG 平移之和。考察 MCG 在构型上的作用。映射类群自然地作用在 $\cB$ 上，即端点落在边界上的简单曲线（相对边界）的非平凡自由同伦类所成之集。回忆这些对应关系：$\cB$ 的元素 $\iff$ 线框（wire frames）$\iff$ 边界裤子 $\cP$。对照图 1，MCG 在 $\cB$ 上的轨道有三种类型，刻画边界裤子 $\cP$ 的位置。
\begin{itemize}[leftmargin=1.6em,itemsep=2pt,topsep=2pt]
\item 从 $\beta_1$ 到 $\beta_1$ 的简单曲线所成的轨道。
  \begin{itemize}[leftmargin=1.4em,itemsep=2pt,topsep=2pt]
  \item $R-\cP$ 连通时只有单个轨道，且 $\operatorname{Sym}(\beta_1,\alpha_1,\alpha_2)=\Z/2\Z$。
  \item $R-\cP$ 断开时是一族轨道。这些轨道由亏格的分割 $g=g_1+g_2$ 与有标号边界 $\{\beta_2,\ldots,\beta_n\}$ 的分割联合分类，并且所得的每个分支的 Euler 示性数都为负。一般地 $\operatorname{Sym}(\beta_1,\alpha_1,\alpha_2)=1$，例外情形是只有一条边界且 $g_1=g_2$。
  \end{itemize}
\item 对每个第二边界的选取各给出一族轨道。特别地，对从 $\beta_1$ 到 $\beta_j$（$j\neq1$）的简单曲线给出一个轨道。所得的曲面 $R-\cP$ 连通，且 $\operatorname{Sym}(\beta_1,\beta_j,\alpha)=1$。
\end{itemize}

考察定理 8，把恒等式
\[
L_1=\sum_{\alpha_1,\alpha_2}\cD(L_1,\ell_{\alpha_1}(R),\ell_{\alpha_2}(R))
+\sum_{j=2}^n\sum_{\alpha}\cR(L_1,L_j,\ell_\alpha(R))
\]
的两端相对于体积 $\dd V$ 在 $R$ 的模空间上积分。对两端作 $L_1$ 的偏导数，以简化量 $\cD$ 与 $\cR$。对右端应用公式 (3)。把右端写成给定轨道类型下各自的求和。对每种轨道类型应用定理 9，即得到用低维模空间体积表示的积分如下。
\begin{itemize}[leftmargin=1.6em,itemsep=2pt,topsep=2pt]
\item 对 $R-\cP$ 连通、从 $\beta_1$ 到 $\beta_1$ 的简单曲线求和时，被加函数是 $\cD$，且
\[
\frac{\partial\cD}{\partial L_1}=H(\ell_{\alpha_1}+\ell_{\alpha_2},L_1),
\]
它是多重曲线长度的函数；所得积分为
\[
\int_0^\infty\int_0^\infty H(x+y,L_1)V_{g-1}(x,y,\widehat{L})\,x y\,\dd x\,\dd y,
\]
其中 $\widehat{L}=(L_2,\ldots,L_n)$。
\end{itemize}
% 注：p.23 末的这个 itemize 只含第一条 bullet，其余两条 bullet 在 p.24（tr03c）开头继续；
% 因文件边界不能切断环境，这里关闭列表，由 tr03c 另开一个新的 itemize 承接。
